% Ad Familiares 3.3 --- headnote
% ad-familiares-03-03, ca. 22--31 May 51 BC, written at Brundisium

\textit{Cicero to Appius Claudius Pulcher, written at Brundisium
toward the end of May 51 BC, on the eve of the sea crossing to
Cilicia. The Perseus dateline places the letter at the end of
May; \S 1 fixes the day of Cicero's arrival at Brundisium as
the eleventh day before the Kalends of June (22 May), and the
letter is plainly written within a few days of that arrival,
while Cicero is waiting on his legate C. Pomptinus before
sailing. This is the third piece in the studied handover
correspondence between two men whom mutual reserve does not
keep from doing necessary business: Cicero is succeeding the
elder brother of his old enemy Clodius in a frontier province,
and is at pains throughout to insist that the transition is
between friends.}

\textit{The substance is two pieces of business raised by Appius's
legate Fabius Vergilianus, who has come down to Brundisium with
messages from his chief. First, the Senate's view --- shared,
Cicero notes, by Appius and reported through Fabius before
Cicero had thought of it --- that reinforcements were needed
for the Cilician command; the consul Servius Sulpicius blocked
the levy, and the Senate's overriding instruction was that
Cicero and Bibulus should set out at once regardless. Second,
the disposition of Appius's existing troops: Appius had written
to the Senate as though he were discharging soldiers wholesale,
but Fabius reports that the discharges had not yet taken place
when he left, and Cicero asks that they not now be carried
through, since the legions are already thin. The tone
throughout is the same careful courtesy as Ad Familiares 3.1
--- the request is framed as confidence rather than complaint,
and the closing line about waiting on Pomptinus before sailing
is meant in part to give Appius time to receive the letter and
respond before the new governor arrives.}
